\section{Постановка задачи}
\label{sec:problem}

Настоящая работа решает задачу построения инфраструктуры автоматической оценки
русскоязычных RAG-систем. Формально задача разбивается на три подзадачи:
(а)~собрать русскоязычный размеченный корпус,
(б)~обучить модель-оценщик,
 (в)~проанализировать её поведение в сравнении с LLM-оценщиком и кросс-язычным базовым классификатором.

\subsection{Формальная постановка задачи классификации}

Пусть $q$ - вопрос пользователя, $D = \{d_1, \dots, d_{|D|}\}$ - упорядоченное 
множество извлечённых ретривером документов, $a$ - ответ генератора. 
Модель-оценщик $M_\theta$ получает на вход объединённую последовательность
\[
x = [q]\,[\text{SEP}]\,[d_1, \dots, d_{|D|}]\,[\text{SEP}]\,[a],
\]
токенизированную в $x = (x_1, \dots, x_T)$, и возвращает три токенных вектора логитов:
\[
M_\theta(x) = \bigl(\ell^{(R)}_t,\, \ell^{(U)}_t,\, \ell^{(A)}_t\bigr)_{t=1}^{T}.
\]

Головы соответствуют трём целевым меткам TRACe:
\begin{itemize}
  \item \textbf{Relevance}: бинарная метка на каждом токене документов $y^{(R)}_t \in \{0, 1\}$. 
  Равна~1, если токен входит в предложение, помеченное LLM-аннотатором как релевантное запросу.
  \item \textbf{Utilization}: бинарная метка на каждом токене документов $y^{(U)}_t \in \{0, 1\}$. 
  Равна~1, если токен входит в предложение, использованное в ответе.
  \item \textbf{Adherence}: бинарная метка на каждом токене ответа $y^{(A)}_t \in \{0, 1\}$. 
  Равна~1, если соответствующее утверждение в ответе поддержано контекстом.
\end{itemize}

Введём две бинарные маски: $m^{(C)}_t$ - маска токенов 
документов (для $y^{(R)}$ и $y^{(U)}$) и $m^{(A)}_t$ - маска токенов ответа (для $y^{(A)}$). 
Тогда функция потерь - взвешенная сумма бинарных кросс-энтропий по валидным токенам:
\[
\mathcal{L}(\theta) = \alpha_R \mathcal{L}_R + \alpha_U \mathcal{L}_U + \alpha_A \mathcal{L}_A,
\]
\[
\mathcal{L}_R = \frac{\sum_t m^{(C)}_t \cdot 
\mathrm{BCE}\bigl(\sigma(\ell^{(R)}_t),\, y^{(R)}_t\bigr)}{\sum_t m^{(C)}_t},
\]
и аналогично для $\mathcal{L}_U$, $\mathcal{L}_A$ (последний - по маске $m^{(A)}_t$).

В работе рассматриваются две схемы весов: 
$(\alpha_R, \alpha_U, \alpha_A) = (1,1,1)$ (базовая) и $(3,3,1)$ (гипотеза о том, что adherence 
- более лёгкая задача и не требует увеличенного веса).

\subsection{Переход от токенных предсказаний к агрегированным метрикам}

При применении модели бинаризация производится по порогам $(\tau_R, \tau_U, \tau_A)$:
\[
\hat{y}^{(R)}_t = \mathbb{1}\bigl[\sigma(\ell^{(R)}_t) \ge \tau_R\bigr].
\]
Для каждой конфигурации пороги подбираются независимо на валидационной
выборке путём сеточного поиска по $\{0{.}1, 0{.}11, \dots, 0{.}9\}$,
максимизируя F1 на уровне токенов (token-level~F1) на соответствующей маске.

Итоговые TRACe-величины восстанавливаются агрегированием токенных предсказаний
 в рамках документа и примера (формулы из раздела~\ref{sec:ragbench}); в работе,
 однако, основной метрикой качества модели-оценщика выступает именно F1 на уровне токенов по каждому
 таргету, так как он не зависит от ошибок, накапливающихся при агрегации.

\subsection{Метрики качества модели-оценщика}

Для сравнения конфигураций используются:
\begin{itemize}
  \item Precision / Recall / F1 на уровне токенов по каждой из трёх голов;
  \item средний F1 по трём таргетам: $\bar{F}_1 = (F_1^{(R)} + F_1^{(U)} + F_1^{(A)}) / 3$ - основной агрегат для ранжирования;
  \item корреляция Пирсона между предсказанными TRACe-скорами на уровне примера и разметкой RAGBench - дополнительная сверка с оригинальной схемой.
\end{itemize}

\subsection{Задачи, решаемые в диссертации}

Исходя из постановки, работа состоит из следующих задач:

\begin{enumerate}
  \item \textbf{LLM-разметка (раздел~\ref{sec:llm_eval}).} Воспроизвести схему разметки 
  из оригинальной статьи на Qwen2.5-72B для русскоязычных примеров, оценить долю формальных ошибок.
  \item \textbf{Русскоязычный корпус (раздел~\ref{sec:dataset}).} Построить RAGBench-RU,
   переведя 7 поддатасетов исходного корпуса, провести контроль качества и опубликовать на HuggingFace.
  \item \textbf{Обучение моделей-оценщиков (разделы~\ref{sec:model}--\ref{sec:grid}).} Обучить модели
   на основе трёх базовых архитектур при нескольких длинах контекста, двух наборах весов функции
   потерь и двух типах головы. Подбор гиперпараметров
   (learning rate для базовой части модели и головы, warmup, dropout) фиксирован для всех запусков.
  \item \textbf{Сравнительный анализ (подразделы~\ref{sec:ablation_weights}--\ref{sec:classifier_vs_llm}).} Оценить влияние длины контекста, весов лосса, архитектуры головы и языка; сопоставить лучший классификатор с LLM-оценщиком.
  \item \textbf{Сервис применения (раздел~\ref{sec:service}).} Реализовать программный интерфейс на 
  базе Streamlit 
  для применения обученных моделей в режимах одиночной и пакетной обработки данных.
\end{enumerate}



